
% =========================================================
\section{Genetic-Algorithm Optimization Framework}
\label{sec:genetic_algorithm}
% =========================================================

\subsection{Optimization problem definition}

The genetic algorithm searches the prescribed design space while retaining the engine
operating point and pre-sizing quantities supplied by the user. The detailed
quasi-one-dimensional, boundary-layer, performance and contour formulations are given
in Sections~\ref{sec:3}--\ref{sec:geometry} and are not repeated here. Within the optimizer, these models are
treated as a coupled candidate-evaluation function

\begin{equation}
\mathcal{S}:\left(\bm{q},\bm{x}_{fixed}\right)
\longrightarrow
C_{F,eff},
\label{eq:coupled_candidate_evaluator}
\end{equation}

where \(\bm{q}\) is the chromosome and \(\bm{x}_{fixed}\) contains the quantities
held constant during one optimization run.

For every candidate, the simulator evaluates the required RocketCEA properties, the
common physical models and the effective thrust coefficient defined in
Eq.~\eqref{eq:effective_thrust_coefficient}. The genetic algorithm therefore operates
on the same physical model used by the simulator rather than on an independent
surrogate.

The Bartz heat-transfer profile is omitted from the repeated fitness calculation
because thermal loading is not part of the current objective. No wall-temperature
input is passed to the evaluator: both boundary-layer closures use the locally derived
adiabatic-wall temperature of Eq.~\eqref{eq:adiabatic_wall_temperature}.

Two evaluation modes are exposed without changing the chromosome or GA operators.
\emph{Quick screening} uses the weak integral closure of
Section~\ref{sec:integral_boundary_layer}; \emph{BLIMP-lite} uses the profile marcher
of Section~\ref{sec:blimp_cebeci_smith}. A solution obtained through Quick screening
is regenerated and reported with BLIMP-lite before it is accepted as the selected
geometry.

% ---------------------------------------------------------
\subsection{Chromosome and fixed quantities}
% ---------------------------------------------------------

\subsubsection{Chromosome definition}

The implemented chromosome contains three real-valued genes:

\begin{equation}
\boxed{
\bm{q}
=
\begin{bmatrix}
\varepsilon & K_b & \theta_{in}
\end{bmatrix}^{T}
},
\label{eq:implemented_chromosome}
\end{equation}

where \(\varepsilon\) is the exit-to-throat area ratio, \(K_b\) is the bell-length
fraction and \(\theta_{in}\) is the initial divergent-wall angle. Their default search
bounds are

\begin{align}
2.0&\leq\varepsilon\leq15.0,
\label{eq:epsilon_gene_bounds}\\
0.60&\leq K_b\leq1.00,
\label{eq:bell_fraction_gene_bounds}\\
20^\circ&\leq\theta_{in}\leq35^\circ.
\label{eq:theta_in_gene_bounds}
\end{align}

All genes are represented as floating-point values. The limits can be changed directly
from the optimization interface.

The exit angle \(\theta_{out}\) is not a gene because it is a derived output of the
contour model. The convergent angle \(\theta_{sub}\) is also excluded from the
chromosome and remains fixed during the search.

\subsubsection{Fixed quantities}

The fixed input vector may be represented as

\begin{equation}
\bm{x}_{fixed}
=
\begin{bmatrix}
p_c & O/F & p_a & r_t & D_c & \alpha & \theta_{sub}
\end{bmatrix}^{T}.
\label{eq:ga_fixed_inputs}
\end{equation}

The default values used by the interface are listed in
Table~\ref{tab:ga_inputs_updated}.

\begin{table}[htbp]
\centering
\caption{Default inputs and their roles during optimization.}
\label{tab:ga_inputs_updated}
\begin{tabular}{llll}
\toprule
Quantity & Symbol & Default value & GA status \\
\midrule
Chamber pressure          & \(p_c\)          & \(30\,\mathrm{bar}\)    & Fixed \\
Mixture ratio             & \(O/F\)          & \(6.5\)                  & Fixed \\
Ambient pressure          & \(p_a\)          & \(1.0\,\mathrm{atm}\)   & Fixed \\
Expansion ratio           & \(\varepsilon\)  & \(5.6\)                  & Optimized \\
Throat radius             & \(r_t\)          & \(0.01728\,\mathrm{m}\) & Fixed \\
Chamber diameter          & \(D_c\)          & \(0.120\,\mathrm{m}\)  & Fixed \\
Reference cone half-angle & \(\alpha\)       & \(15^\circ\)            & Fixed \\
Bell-length fraction      & \(K_b\)          & \(0.80\)                 & Optimized \\
Initial divergent angle   & \(\theta_{in}\)  & \(30^\circ\)            & Optimized \\
Convergent angle          & \(\theta_{sub}\) & \(50^\circ\)            & Fixed \\
\bottomrule
\end{tabular}
\end{table}

The recovery factor is not included because it is derived locally from the
interpolated RocketCEA Prandtl number rather than supplied as an independent input.
Since \(\varepsilon\) is
part of the chromosome, the exit properties returned by RocketCEA are candidate
dependent and cannot be treated as a single common result for the complete population.

% ---------------------------------------------------------
\subsection{Fitness-function definition}
% ---------------------------------------------------------

The scalar objective maximized by the genetic algorithm is the effective ambient
thrust coefficient introduced in Section~\ref{sec:loss_corrected_model}:

\begin{equation}
\boxed{
J(\bm{q})
=
C_{F,eff}(\bm{q})
}.
\label{eq:implemented_ga_fitness}
\end{equation}

Explicitly,

\begin{equation}
J(\bm{q})
=
\eta_{div}(\bm{q})\eta_{BL}(\bm{q})
C_{F,mom}^{CEA}(\bm{q})
+
\varepsilon(\bm{q})
\frac{p_e(\bm{q})-p_a}{p_c}
-C_{F,fric}(\bm{q}).
\label{eq:expanded_ga_fitness}
\end{equation}

The objective is maximized directly, so no sign inversion is required. It contains a
loss-corrected momentum contribution, the signed ambient pressure-thrust contribution
and integrated wall friction. It does not include additional empirical, mass,
manufacturing-cost or thermal-limit penalties.

% ---------------------------------------------------------
\subsection{Candidate validation and constraint handling}
% ---------------------------------------------------------

The gene bounds in Eqs.~\eqref{eq:epsilon_gene_bounds}--%
\eqref{eq:theta_in_gene_bounds} provide the first level of constraint enforcement.
Each generated candidate is then passed to the common simulator validation and to the
models defined in Sections~\ref{sec:3}--\ref{sec:geometry}. The geometric
admissibility conditions are part of the contour model of
Section~\ref{sec:geometry} and are not duplicated here.

A candidate is rejected if the common simulator reports an invalid input, an invalid
contour, a singular numerical system, a non-finite result or a non-positive effective
thrust coefficient. Candidates classified as separated by RocketCEA are also rejected.

Constraint handling uses a death-penalty strategy:

\begin{equation}
J_{GA}(\bm{q})
=
\begin{cases}
C_{F,eff}(\bm{q}),
&\text{valid, finite and attached-flow candidate},\\[1.5mm]
0,
&\text{otherwise}.
\end{cases}
\label{eq:ga_death_penalty}
\end{equation}

There is no continuous penalty proportional to the magnitude of a violation.

% ---------------------------------------------------------
\subsection{Genetic-algorithm configuration}
% ---------------------------------------------------------

The optimizer is implemented using PyGAD \cite{PyGAD}. Its default configuration is summarized in
Table~\ref{tab:ga_configuration_updated}.

\begin{table}[htbp]
\centering
\caption{Default configuration of the implemented genetic algorithm.}
\label{tab:ga_configuration_updated}
\begin{tabular}{ll}
\toprule
Parameter & Implemented value \\
\midrule
Chromosome representation & Three real-valued genes \\
Population size & \(N_p=100\) individuals \\
Maximum generations & \(G_{max}=300\) \\
Mating parents & 20 \\
Parent selection & Tournament selection \\
Crossover & Uniform crossover \\
Crossover probability & \(P_c=0.85\) \\
Mutation & Adaptive mutation \\
Adaptive mutation percentages & \([67\%,34\%]\) \\
Elitism & Best three individuals \\
Saturation limit & 40 generations \\
Default evaluation mode & Four worker processes \\
Repeated-candidate cache & Enabled \\
Saved full solution history & No \\
Explicit random seed & No \\
\bottomrule
\end{tabular}
\end{table}

The settings are user-configurable. The implemented validation requires

\begin{equation}
N_p\geq4,
\qquad
G_{max}\geq1,
\qquad
2\leq N_{parents}\leq N_p.
\label{eq:minimum_ga_settings}
\end{equation}

\subsubsection{Tournament selection}

Tournament selection gives preference to candidates with higher
\(C_{F,eff}\) while preserving stochastic exploration. Candidate parents are sampled
from the current population and the best member of each tournament is selected for
reproduction. With the default settings, 20 parents participate in each generation.

\subsubsection{Uniform crossover}

Uniform crossover constructs each offspring by selecting each gene from either
parent. For example, given

\begin{align}
\bm{q}^{(1)}
&=
\begin{bmatrix}
\varepsilon^{(1)} & K_b^{(1)} & \theta_{in}^{(1)}
\end{bmatrix}^{T},\\
\bm{q}^{(2)}
&=
\begin{bmatrix}
\varepsilon^{(2)} & K_b^{(2)} & \theta_{in}^{(2)}
\end{bmatrix}^{T},
\end{align}

one possible offspring is

\begin{equation}
\bm{q}^{(o)}
=
\begin{bmatrix}
\varepsilon^{(1)} & K_b^{(2)} & \theta_{in}^{(1)}
\end{bmatrix}^{T}.
\label{eq:uniform_crossover_example}
\end{equation}

Crossover is applied with probability \(P_c=0.85\).

\subsubsection{Adaptive mutation}

Adaptive mutation changes the mutation intensity according to candidate fitness. The
configured percentages are

\begin{equation}
\bm{p}_m
=
\begin{bmatrix}
67\% & 34\%
\end{bmatrix}.
\label{eq:adaptive_mutation_percentages}
\end{equation}

Lower-fitness candidates receive the larger mutation percentage, promoting
exploration, while higher-fitness candidates receive the smaller percentage,
promoting local refinement. For the three-gene chromosome, these percentages
correspond approximately to mutating two genes and one gene, respectively. All
mutations remain bounded by the defined gene spaces.

\subsubsection{Elitism}

The best three candidates are transferred directly to the following generation,

\begin{equation}
N_{\mathrm{elite}}=3.
\label{eq:ga_elitism}
\end{equation}

These individuals are preserved without being subjected to crossover or mutation, ensuring that the highest-fitness nozzle geometries already identified by the optimizer are not lost through the stochastic genetic operators. Since only three individuals are retained from a population of 100, elitism preserves the best solutions while leaving most of the population available for continued exploration of the design space. 

% ---------------------------------------------------------
\subsection{Evaluation acceleration}
% ---------------------------------------------------------

The optimizer supports serial, threaded and multiprocess evaluation. The default mode
uses a persistent pool of four worker processes. Each process receives an isolated
RocketCEA working directory, preventing simultaneous Fortran evaluations from using
the same temporary files.

An exact-match fitness cache stores previously evaluated chromosomes. When an
identical chromosome reappears, its stored fitness is reused. The ideal expansion
ratio associated with the fixed values of \(p_c\), \(O/F\) and \(p_a\) is cached
separately because it does not depend on the candidate chromosome.

These measures reduce evaluation time without changing population generation,
selection, crossover, mutation, elitism or stopping logic.

% ---------------------------------------------------------
\subsection{Stopping, cancellation and reported progress}
% ---------------------------------------------------------

The optimization terminates when the first of the following conditions is met:

\begin{enumerate}
    \item the maximum number of generations is reached;
    \item the best fitness does not improve for the selected saturation interval,
    equal to 40 generations by default;
    \item the user requests cancellation through the interface.
\end{enumerate}

Cancellation is checked at the end of each generation. The current generation is
therefore allowed to finish before execution stops.

During execution, the interface displays the completed generation, percentage,
elapsed time, estimated remaining time and generation rate. It also reports the best
\(C_{F,eff}\) and the current chromosome. If no candidate with positive fitness is
found, the optimization is declared unsuccessful.

Otherwise, the returned result is

\begin{equation}
\bm{y}_{GA}
=
\begin{bmatrix}
\varepsilon^* &
K_b^* &
\theta_{in}^* &
\theta_{out}^* &
p_e^* &
C_{F,eff}^* &
N_g
\end{bmatrix}^{T},
\label{eq:ga_output}
\end{equation}

where \(N_g\) is the number of completed generations. The ambient operating-mode
classification is also retained. Because no explicit random seed is assigned,
repeated runs with identical inputs may return slightly different chromosomes.


\subsection{Algorithmic limitations and possible extensions}
% ---------------------------------------------------------

The implemented GA solves a single-objective, single-operating-point problem. It is
stochastic and does not guarantee the global optimum. Its result also depends on the
accuracy and validity range of the coupled simulator used to evaluate each chromosome.

Future algorithmic extensions may include multi-objective optimization, robustness
across several ambient pressures, explicit uncertainty propagation, repeated runs
with controlled random seeds, convergence-statistics reporting and surrogate-assisted
evaluation. Physical or geometric extensions should be introduced in their respective
model sections and then exposed to the GA only through the chromosome, constraints or
candidate-evaluation function.

% ---------------------------------------------------------
\subsection{Optimization workflow}
% ---------------------------------------------------------

Figure~\ref{fig:ga_workflow_updated} summarizes the implemented algorithm without
repeating the internal formulations of the physical simulator.

\begin{figure}[htbp]
\centering
\begin{tikzpicture}[
    node distance=7mm and 18mm,
    every node/.style={font=\small},
    process/.style={
        rectangle,
        rounded corners,
        draw,
        align=center,
        minimum width=45mm,
        minimum height=8mm
    },
    decision/.style={
        diamond,
        draw,
        aspect=2.1,
        align=center,
        inner sep=1.5pt
    },
    terminal/.style={
        ellipse,
        draw,
        align=center,
        minimum width=35mm,
        minimum height=8mm
    },
    arrow/.style={
        -{Latex[length=2mm]},
        thick
    }
]

% =========================================================
% MAIN COLUMN
% =========================================================

\node[terminal] (start) {
Start optimization
};

\node[process, below=of start] (settings) {
Read fixed inputs, gene bounds\\
and GA configuration
};

\node[process, below=of settings] (population) {
Generate population\\
\(\bm{q}=[\varepsilon,K_b,\theta_{in}]^T\)
};

\node[process, below=of population] (evaluate) {
Evaluate candidates using the\\
coupled simulator
};

\node[decision, below=of evaluate] (valid) {
Valid, finite and\\
attached solution?
};

\node[process, below=of valid] (fitness) {
Assign\\
\(J=C_{F,\mathrm{eff}}\)
};

\node[decision, below=of fitness] (all) {
All candidates\\
evaluated?
};

\node[process, below=of all] (progress) {
Update best solution, progress\\
and preserve elites
};

\node[decision, below=of progress] (stop) {
Stopping criterion\\
satisfied?
};

\node[terminal, below=of stop] (output) {
Return best chromosome and performance
};


% =========================================================
% RIGHT-HAND NODES
% =========================================================

\node[process, right=22mm of valid] (zero) {
Assign \(J=0\)
};

\node[process, right=22mm of all] (next) {
Evaluate next\\
candidate
};

\node[process, right=22mm of stop] (reproduce) {
Tournament selection, uniform\\
crossover and adaptive mutation
};


% =========================================================
% MAIN FLOW
% =========================================================

\draw[arrow] (start) -- (settings);
\draw[arrow] (settings) -- (population);
\draw[arrow] (population) -- (evaluate);
\draw[arrow] (evaluate) -- (valid);

% Validity decision
\draw[arrow]
    (valid) -- node[left]{Yes} (fitness);

\draw[arrow]
    (valid) -- node[above]{No} (zero);

\draw[arrow]
    (fitness) -- (all);


% =========================================================
% INVALID CANDIDATE:
% DIRECTLY REJOINS THE "ALL CANDIDATES?" DECISION
% =========================================================

\draw[arrow]
    (zero.south)
    to[out=-105,in=25,looseness=1.05]
    (all.25);


% =========================================================
% CANDIDATE LOOP
% =========================================================

\draw[arrow]
    (all) -- node[left]{Yes} (progress);

\draw[arrow]
    (all) -- node[above]{No} (next);

% Next candidate returns to evaluator
\draw[thick]
    (next.east)
    -- (10.15,-11.58);


% =========================================================
% GENERATION LOOP
% =========================================================

\draw[arrow]
    (progress) -- (stop);

\draw[arrow]
    (stop) -- node[left]{Yes} (output);

\draw[arrow]
    (stop) -- node[above]{No} (reproduce);

% New generation returns to evaluator
\draw[arrow]
    (reproduce.east)
    -- ++(10mm,0)
    |- ([xshift=8mm]evaluate.east)
    -- (evaluate.east);

\end{tikzpicture}

\caption{Flowchart of the implemented genetic-algorithm optimization.}
\label{fig:ga_workflow_updated}
\end{figure}

% ---------------------------------------------------------

